% !TEX program = pdflatex
% The commutator structure of canonical transcorrelation: the ansatz, the
% expansion of the sheets, and the full list of commutators the derivation
% needs, with the verified rule for each.
\documentclass[11pt]{article}
\usepackage[T1]{fontenc}
\usepackage[utf8]{inputenc}
\usepackage{lmodern}
\usepackage{amsmath,amssymb}
\usepackage[margin=1in]{geometry}
\usepackage{booktabs}
\usepackage[hidelinks]{hyperref}

\newcommand{\E}[2]{\hat E^{#1}_{#2}}
\newcommand{\ad}{\hat a^\dagger}
\newcommand{\hc}{\mathrm{h.c.}}

\title{\bfseries The commutator structure of canonical transcorrelation}
\author{}
\date{}

\begin{document}
\maketitle
\vspace{-2.5em}

\noindent
This is the skeleton of the CT derivation as it is written on the sheets: the
ansatz, the second-order expansion, and the decomposition of the commutators
that has to be carried out. Section~4 lists every block the derivation needs,
with the rule that produces it and its present status. Nothing here is a new
result; the point is to fix the structure once, so the block-by-block algebra
of \texttt{progress\_note/singles\_derivation\_plan.md} has a single home.

\section*{1. The ansatz}

\begin{align}
  \hat H&=\hat F+\hat V,
   & \hat F&=f^q_p\,\E{q}{p},
   & \hat V&=\tfrac12\,g^{rs}_{pq}\,\E{pq}{rs},
   \label{eq:split}\\[2pt]
  \hat A&=\hat T-\hat T^\dagger,
   & \hat T&=\hat T_1+\hat T_2,
   & \hat T_1&=t^a_i\,\E{i}{a},
   & \hat T_2&=\tfrac12\,t^{ab}_{ij}\,\E{ij}{ab}.
   \label{eq:ansatz}
\end{align}
The transformation is
\begin{equation}
  \hat{\bar H}=e^{-\hat A}\,\hat H\,e^{\hat A}
   =e^{\hat A^\dagger}\hat H\,e^{\hat A},
  \label{eq:transform}
\end{equation}
the two forms being the same because $A$ is anti-Hermitian,
$\hat A^\dagger=-\hat A$. Expanding and keeping terms through second order,
\begin{equation}
  \hat{\bar H}=\hat H+[\hat H,\hat A]
   +\tfrac12\bigl[[\hat F,\hat A],\hat A\bigr]+\cdots,
  \label{eq:second}
\end{equation}
where the Fock operator replaces $\hat H$ in the double commutator --- the
second-order truncation, not an identity. In the F12 case this is Yanai and
Shiozaki's own truncation
\cite{Yanai:2012cj},
\[
  \hat{\bar H}^{F_{12}}
   \approx\hat H+[\hat H,\hat A_{F_{12}}]_{1,2}
   +\tfrac12\bigl[[\hat F,\hat A_{F_{12}}]_{1,2},\hat A_{F_{12}}\bigr]_{1,2},
\]
and the same form with the explicit single excitation added is
Kumar \emph{et al.}\ \cite{Kumar2022}, eqs.~(3)--(4). Two conventions have to be
stated with it:

\begin{itemize}
  \item \textbf{The $[\dots]_{1,2}$ subscript.} Only one- and two-body
        operators generated by a commutator are kept; the three-body pieces are
        folded back with the extended normal ordering (cumulants) of
        Kutzelnigg and Mukherjee \cite{Kutzelnigg:1997}. This is an
        approximation and is recorded as one.
  \item \textbf{The single excitation.} $\hat T_1$ is Kumar's addition;
        Yanai and Shiozaki's $\hat A^{F12}$ contains the projected geminal
        only. Everything below is written with $\hat T_1$ present and reduces
        to the geminal-only case at $t^a_i=0$.
\end{itemize}

\section*{2. The decomposition}

Because $\hat A=\hat T-\hat T^\dagger$ and $\hat H$ is Hermitian,
\begin{equation}
  [\hat H,\hat A]=[\hat H,\hat T]+[\hat H,\hat T]^\dagger,
  \qquad
  \bigl([\hat H,\hat T]\bigr)^\dagger=-[\hat H,\hat T^\dagger],
  \label{eq:hc}
\end{equation}
so only the $\hat T$ half of every commutator has to be computed. Splitting
$\hat H=\hat F+\hat V$ and $\hat T=\hat T_1+\hat T_2$,
\begin{align}
  [\hat H,\hat T]&=[\hat F,\hat T_1]+[\hat F,\hat T_2]
                  +[\hat V,\hat T_1]+[\hat V,\hat T_2],
   \label{eq:first}\\[2pt]
  [\hat F,\hat A]&=\hat C=\hat C_1+\hat C_2,
   \qquad
  \hat C_i=[\hat F,\hat A_i]=[\hat F,\hat T_i]+[\hat F,\hat T_i]^\dagger .
   \label{eq:C}
\end{align}
The double commutator collapses before it is expanded. Since $\hat A$ is
anti-Hermitian, $\hat C=[\hat F,\hat A]$ is Hermitian, so
$[\hat C,\hat T]^\dagger=-[\hat C,\hat T^\dagger]$ and, using \eqref{eq:C},
\begin{equation}
  \tfrac12\bigl[[\hat F,\hat A],\hat A\bigr]
   =\tfrac12[\hat C,\hat T]+\hc
   =\tfrac12[\hat C_1,\hat T]+\tfrac12[\hat C_2,\hat T]+\hc .
  \label{eq:double}
\end{equation}
Expanding the outer commutator in the two $\hat T$'s,
\begin{equation}
  \tfrac12\bigl[[\hat F,\hat A],\hat A\bigr]
   =\tfrac12\Bigl(
      [\hat C_1,\hat T_1]+[\hat C_1,\hat T_2]
     +[\hat C_2,\hat T_1]+[\hat C_2,\hat T_2]\Bigr)+\hc,
   \label{eq:doubleexpanded}
\end{equation}
which is the statement of the sheets at steps \textcircled{5} and
\textcircled{6}: four double commutators, eight commutators in the $\hat T$'s.

Two remarks on \eqref{eq:double}. The inner commutator is with $\hat F$ --- it
is what produces $\hat C$ --- and only the \emph{outer} one is with $\hat T$;
writing the middle step with $\hat A$ inside,
$\tfrac12[[\hat C_i,\hat A],\hat T]$, gives a different operator, because
$[[\hat C_i,\hat A],\hat T]=[[\hat C_i,\hat T],\hat T]
-[[\hat C_i,\hat T^\dagger],\hat T]$. And the conjugate half is not an
independent computation: \eqref{eq:double} says the whole double commutator is
$[\hat C,\hat T]$ plus its own conjugate, so only four of the eight
commutators of \eqref{eq:doubleexpanded} have to be evaluated.
Both readings are checked on explicit matrices in
\path{tests/test_ct_structure.py}: \eqref{eq:double} holds to machine
precision, and the reading with $\hat A$ inside is off by a matrix norm of
order $10^4$.

The four blocks are \emph{not} pairwise equal ---
\begin{equation}
  \hat T_1\hat T_2-\hat T_2\hat T_1\neq0
\end{equation}
(4 terms), so the Jacobi shortcut
$[[\hat F,\hat T_1],\hat T_2]=[[\hat F,\hat T_2],\hat T_1]$ is invalid and the
difference $-[[\hat T_1,\hat T_2],\hat F]$ has to be carried. Any place that
used that shortcut has to be redone.

\section*{3. The rules the blocks are built from}

Every block in \S4 is a sum of products of two $\hat E$ operators, so the whole
derivation needs only the four product rules and the four commutators of the
$\hat E$ operators, at ranks $(1,1)$, $(1,2)$, $(2,1)$ and $(2,2)$. In the
paired order --- every join coefficient $+1$, the residual's blocks listed
re-paired first, then the left leftovers, then the right leftovers --- they are
\begin{align}
  \bigl[\E{q}{p},\E{s}{r}\bigr]
   &=\delta^s_p\,\E{q}{r}-\delta^q_r\,\E{s}{p},
   \label{eq:c11}\\[2pt]
  \bigl[\E{q}{p},\E{tu}{rs}\bigr]
   &=\delta^t_p\,\E{qu}{rs}+\delta^u_p\,\E{qt}{sr}
     -\delta^q_r\,\E{tu}{ps}-\delta^q_s\,\E{ut}{pr},
   \label{eq:c12}\\[2pt]
  \bigl[\E{tu}{rs},\E{q}{p}\bigr]
   &=-\bigl[\E{q}{p},\E{tu}{rs}\bigr],
   \label{eq:c21}
\end{align}
and the $(2,2)$ commutator, whose four-body part cancels and which leaves eight
three-body and four two-body terms,
\begin{align}
  \bigl[\E{rs}{pq},\E{uw}{tv}\bigr]
   &=\delta^u_p\,\E{rsw}{tqv}+\delta^w_p\,\E{rsu}{vqt}
     +\delta^u_q\,\E{srw}{tpv}+\delta^w_q\,\E{sru}{vpt}\notag\\
   &\quad+\delta^u_p\delta^w_q\,\E{rs}{tv}
         +\delta^w_p\delta^u_q\,\E{rs}{vt}\notag\\
   &\quad-\delta^r_t\,\E{uws}{pvq}-\delta^s_t\,\E{uwr}{qvp}
         -\delta^r_v\,\E{wus}{ptq}-\delta^s_v\,\E{wur}{qtp}\notag\\
   &\quad-\delta^r_t\delta^s_v\,\E{uw}{pq}
         -\delta^s_t\delta^r_v\,\E{uw}{qp}.
   \label{eq:c22}
\end{align}
All four are verified against explicit Jordan--Wigner matrices over every label
assignment, in the spin-orbital \emph{and} the spin-free Fock space, by
\path{check_paired_commutators.py} in the singles directory: 0 failing
assignments, worst deviation $0.0$.

\section*{4. The blocks}

\begin{center}
\footnotesize
\begin{tabular}{@{}p{4.4cm}p{3.2cm}p{1.6cm}p{1.5cm}p{3.6cm}@{}}
\toprule
block & rule used & rank out & terms & status \\
\midrule
$[\hat F,\hat T_1]$ & \eqref{eq:c11} & 1 & 2 & verified \\
$[\hat F,\hat T_2]$ & \eqref{eq:c12} & 2 & 4 & verified \\
$[\hat V,\hat T_1]$ & \eqref{eq:c12} & 2 & 4 & verified \\
$[\hat V,\hat T_2]$ & \eqref{eq:c22} & 2 (+3) & 10 & \emph{blocked} \\
$[\hat T_1,\hat T_2]$ & \eqref{eq:c12} & 2 & 4 & verified \\
\midrule
$[\hat F,\hat A_1]$ & \eqref{eq:c11}, \eqref{eq:hc} & 1 & 2 & verified \\
$[\hat F,\hat A_2]$ & \eqref{eq:c12}, \eqref{eq:hc} & 2 & 4 & verified \\
$\tfrac12[[\hat F,\hat A_1],\hat T_1]$ & \eqref{eq:c11} & $\le2$ & 2 & verified \\
$\tfrac12[[\hat F,\hat A_1],\hat T_2]$ & \eqref{eq:c12} & $\le2$ & 8 & verified \\
$\tfrac12[[\hat F,\hat A_2],\hat T_1]$ & \eqref{eq:c12} & $\le2$ & 8 & \emph{blocked} \\
$\tfrac12[[\hat F,\hat A_2],\hat T_2]$ & \eqref{eq:c22} & $\le2$ & 80 & \emph{blocked} \\
\bottomrule
\end{tabular}
\end{center}

\noindent
\textbf{The three blocked blocks are blocked by one bug, not by the algebra.}
The blocks that fail their numerical check are exactly those whose expansion
needs the $(2,2)$ rule \eqref{eq:c22} --- which is itself verified. The fault is
in the canonicalisation step of the symbolic engine
(\texttt{ctformula/algebra.py}, \texttt{canon\_op}): it brings an operator to
canonical form by sorting the upper and the lower index row \emph{independently}
and taking the sign of the two sorts. For a spin-orbital monomial that is
correct --- the monomial is fixed by the two multisets, up to sign. For a
spin-free $\hat E$ it is not: sorting the two rows independently changes the
\emph{pairing} of upper to lower, which is part of the operator, so terms that
differ only by such a change are wrongly merged and their coefficients
inherited from a relisting sign that the spin-free operator does not have. The
engine's own spin-orbital checks pass, and its spin-free ones fail, on exactly
the same expressions --- which is the signature of this bug.

The fix is to canonicalise a spin-free $\hat E$ as a set of pairs $(u_i,l_i)$
with the sign of the pair permutation, keeping the row-wise canonical form for
the spin-orbital monomials; the convention must be carried through the engine,
because it cannot be read off the operator alone.
Until that is done, the three blocked blocks must not be quoted: the term lists
they produce now carry the wrong pairing on some terms, and the failing numeric
sweeps are what exposes it.

\section*{5. What is verified, and what is not}

\begin{itemize}
  \item \textbf{Verified.} The four rules of \S3, in both conventions, over
        every label assignment; and the eight blocks marked verified in \S4, by
        a Fock-space sweep and an independent symbolic re-derivation, which is
        what \path{tests/test_ct.py} reports (3 of its 4 checks pass).
  \item \textbf{Blocked.} $[\hat V,\hat T_2]$,
        $\tfrac12[[\hat F,\hat A_2],\hat T_1]$ and
        $\tfrac12[[\hat F,\hat A_2],\hat T_2]$ --- and with them the explicit
        $\bar h$ and $\bar g$, which cannot be collected until they are right.
  \item \textbf{Open beyond that.} The fold of the three-body pieces back into
        one- and two-body form ($[\dots]_{1,2}$), which is an approximation and
        has to be stated as one wherever the result is used; and the
        multireference generalisation, which needs the density and cumulant
        terms of the general reference rather than the vacuum contractions used
        here.
\end{itemize}

\begin{thebibliography}{9}
\bibitem{Kutzelnigg:1997} W.~Kutzelnigg and D.~Mukherjee,
  \emph{Normal order and extended Wick theorem for a multiconfiguration
  reference wave function}, J.~Chem.~Phys.\ \textbf{107}, 432 (1997).
\bibitem{Yanai:2012cj} T.~Yanai and T.~Shiozaki,
  \emph{Canonical transcorrelated theory with projected Slater-type geminals},
  J.~Chem.~Phys.\ \textbf{136}, 084107 (2012).
\bibitem{Kumar2022} A.~Kumar, A.~Asthana, C.~Masteran, E.~F.~Valeev,
  Y.~Zhang, L.~Cincio, S.~Tretiak and P.~A.~Dub,
  \emph{Quantum simulation of molecular electronic states with a
  transcorrelated Hamiltonian}, J.~Chem.~Theory Comput.\ \textbf{18}, 5312
  (2022).
\end{thebibliography}

\end{document}
